Source : calculs des auteurs : modèle de foyers calé sur les statistiques de l'impôt sur le revenu 2024 (DGFiP) ; coûts des niches : Évaluation des voies et moyens PLF 2026, CPO (2026) ; foncier : DVF et comptes de patrimoine (Insee) ; impôt sur les sociétés : recettes 2025 (Insee) ; annexes B à F.\par Recettes de 2032, en Md€ de 2025 (montants rapportés à l'économie de 2025), arrondies au milliard. Gris : total de chaque bloc. Impôt sur le revenu : décomposition séquentielle avant réaction des comportements (l'ordre change la taille de chaque ligne, pas le total). Taxe foncière communale assise sur la valeur de marché, en deux parts (sol et bâti), à recette communale inchangée : le total ne change pas, mais la charge se redistribue entre logements. Taxe nationale sur le sol : 15 Md€ bruts, montant affiché net de sa déductibilité des bénéfices. Impôt sur les sociétés : rendements prudents, après réaction des entreprises. En 2032, la déduction immédiate coûte cher parce que l'État rembourse aussi la transition (la marge de l'ensemble tombe à −8 Md€ en 2033) ; une fois celle-ci passée, vers 2040, la marge de l'ensemble atteint environ 5 Md€. Niches : coûts des Voies et moyens du PLF 2026, hors dispositifs qui s'éteignent avant 2032 en droit actuel (Girardin). Quotients conjugal et familial : la Cour des comptes les chiffrait à 11,1 et 16,5 Md€ en 2021. Fin de la CSG déductible : 16 Md€ si la mesure était prise isolément. Hachures : options non retenues, non comptées dans le solde (coût après réaction des contribuables ; entre parenthèses, le solde qui en résulterait). Comptés dans la partie III du rapport et retirés ici : recentrage du crédit d'impôt recherche, mécénat, zonages et crédit famille (mesure 2), fin de l'effacement des plus-values à la transmission (mesure 3).
